\begin{center}\begin{tikzpicture}[font=\small,>=stealth,line join=round]
\draw[gray!40,rounded corners=3pt] (0,0) rectangle (14.2,8.3);\node at (7.1,7.700000000000001) {为什么可以说顺时针？准确地看共同点处的切线};
\begin{scope}[shift={(4,3.9)},x=1.25cm,y=1.25cm]\draw[->,gray] (-2.2,0)--(2.2,0) node[right]{$x$};\draw[->,gray] (0,-1.3)--(0,2.4) node[above]{$y$};\draw[violet,thick,solid] (-.8,-0.10903548889591241)--(.8,2.1090354888959126);\draw[blue!75!black,thick,dashed] (-.8,0.4454822555520438)--(.8,1.5545177444479563);\draw[orange!80!black,thick,densely dotted] (-.8,1.5545177444479563)--(.8,0.4454822555520438);\draw[teal,thick,dashdotted] (-.8,2.1090354888959126)--(.8,-0.10903548889591241);\fill[orange!80!black] (0,1) circle(2pt);\draw[->,black] (1.15,2.15) arc[start angle=45,end angle=-45,radius=1.62635];\end{scope}\node[anchor=west,align=left,text width=5.1cm] at (8.2,4.2) {共同点$(0,1)$处的切线\\[8pt]底数$4\to2\to\frac12\to\frac14$\\[8pt]方向顺时针改变\\[8pt]高中导数核对：\\[5pt]$f\,\!\prime(0)=\ln a$};
\node[text width=12.6cm,align=center] at (7.1,.65) {底数减小，切线斜率随之减小；a=1时为水平线y=1，通常单列为常数函数。};
\end{tikzpicture}\end{center}
